% ============================================================
%  Dízimas Periódicas — Apostila Premium | PratiqueJá
%  Engine: XeLaTeX
% ============================================================
\documentclass[12pt]{article}
\usepackage{../../pratiqueja}

% \hlm — destaque de resultado em modo-math (cor sem bold)
\newcommand{\hlm}[1]{{\color{darkblue}#1}}

\setsubject{Dízimas Periódicas}

\begin{document}
\pagenumbering{arabic}
\setstretch{1.2}
\color{bodytext}

\heroheader{Dízimas Periódicas}%
           {Fração Geratriz e o Método Algébrico}%
           {Matemática · Intermediário}

\setlength{\columnsep}{18pt}
\setlength{\columnseprule}{0.5pt}
\def\columnseprulecolor{\color{exborder}}

\begin{multicols}{2}

%─────────────────────────────────────────────────────────────
\TW{Def}{badgepurple}{Definição}{Dízima periódica e números racionais}

\regra{Uma \textbf{dízima periódica} é um decimal cuja parte decimal
se repete infinitamente. A parte que se repete é o \textbf{período},
indicado por $\overline{\;\cdot\;}$.}

\small\color{bodytext}%
Toda dízima periódica é um \textbf{número racional}: pode ser escrita
como $\frac{a}{b}$. Essa fração é a \textbf{fração geratriz}.

\exemp{%
  $0{,}333\ldots = 0{,}\overline{3}$\\[2pt]
  $1{,}272727\ldots = 1{,}\overline{27}$\\[2pt]
  $0{,}1666\ldots = 0{,}1\overline{6}$%
}

\SH{Tipos}
\obs{\textbf{Simples:} a parte decimal contém \emph{só} o período.
Ex.: $0{,}\overline{6}$, $2{,}\overline{56}$.}
\obs{\textbf{Composta:} há uma parte \emph{não periódica} antes do
período. Ex.: $0{,}5\overline{2}$, $1{,}68\overline{74}$.}

%─────────────────────────────────────────────────────────────
\TW{Alg}{darkblue}{Método Algébrico}{De onde vêm as fórmulas}

\regra{Chame a dízima de $x$. Multiplique por potências de $10$ para
alinhar os períodos e \textbf{subtraia}: a parte periódica se cancela
e sobra uma equação inteira.}

\SH{Simples — $0{,}\overline{7}$}
\exemp{%
  $x = 0{,}777\ldots$\\
  $10x = 7{,}777\ldots$\\
  $10x - x = 7 \;\Rightarrow\; 9x = 7$\\
  $x = \hlm{\dfrac{7}{9}}$%
}
\nota{O denominador tem $p$ noves porque
$10^p x - x = (10^p - 1)x$.}

\SH{Composta — $1{,}5\overline{3}$}
\exemp{%
  $x = 1{,}5333\ldots$\\
  $10x = 15{,}333\ldots$ {\footnotesize(desloca não-periódica)}\\
  $100x = 153{,}333\ldots$ {\footnotesize(desloca período)}\\
  $100x - 10x = 138 \;\Rightarrow\; 90x = 138$\\
  $x = \dfrac{138}{90} = \hlm{\dfrac{23}{15}}$%
}
\nota{$10^{n+p} - 10^n$ = $p$ noves seguidos de $n$ zeros.}

%─────────────────────────────────────────────────────────────
\TW{Dica}{badgegreen}{Atalho Mental}{Cálculo rápido da geratriz}

\regra{\textbf{Simples:} numerador $=$ período; denominador $=$ $p$
noves.\\[3pt]
\textbf{Composta:} numerador $=$ (número sem vírgula) $-$ (parte não
periódica); denominador $=$ $p$ noves seguidos de $n$ zeros.}

\SH{Simples}
\exemp{%
  $0{,}\overline{3} = \dfrac{3}{9} = \hlm{\dfrac{1}{3}}$\\[3pt]
  $0{,}\overline{37} = \hlm{\dfrac{37}{99}}$\\[3pt]
  $0{,}\overline{259} = \dfrac{259}{999} = \hlm{\dfrac{7}{27}}$\\[3pt]
  $2{,}\overline{56} = 2 + \dfrac{56}{99} = \hlm{\dfrac{254}{99}}$%
}

\SH{Composta}
\exemp{%
  $0{,}5\overline{2} = \dfrac{52 - 5}{90} = \hlm{\dfrac{47}{90}}$\\[3pt]
  $0{,}1\overline{6} = \dfrac{16 - 1}{90} = \dfrac{15}{90} = \hlm{\dfrac{1}{6}}$\\[3pt]
  $1{,}68\overline{74} = \dfrac{16874 - 168}{9900} = \hlm{\dfrac{8353}{4950}}$%
}

%─────────────────────────────────────────────────────────────
\TW{FG}{darkblue}{Fração Geratriz — Simples}{Fórmula e exemplos}

\regra{Período de $p$ algarismos; o denominador $10^p - 1$ é sempre
$p$ noves $(9, 99, 999\ldots)$:}
\formula{$x = \dfrac{x\cdot10^{p} - x}{10^{p} - 1}$}

\SH{$0{,}\overline{7}$ \quad $(p=1)$}
\exemp{%
  $\dfrac{7{,}\overline{7} - 0{,}\overline{7}}{10 - 1}
   = \hlm{\dfrac{7}{9}}$ \;{\footnotesize(irredutível)}%
}

\SH{$0{,}\overline{259}$ \quad $(p=3)$}
\exemp{%
  $\dfrac{259{,}\overline{259} - 0{,}\overline{259}}{999}
   = \dfrac{259}{999}$\\
  MDC$(259,999) = 37 \;\Rightarrow\; \hlm{\dfrac{7}{27}}$%
}

\needspace{6\baselineskip}
\SH{$2{,}\overline{56}$ \quad $(p=2)$}
\exemp{%
  $\dfrac{256{,}\overline{56} - 2{,}\overline{56}}{99}
   = \dfrac{254}{99}$\\
  MDC$(254,99) = 1 \;\Rightarrow\; \hlm{\dfrac{254}{99}}$%
}

%─────────────────────────────────────────────────────────────
\TW{FG}{badgepurple}{Fração Geratriz — Composta}{Com parte não periódica}

\regra{$n$ = dígitos \textbf{não periódicos}; $p$ = dígitos
\textbf{periódicos}. Denominador $10^{n+p} - 10^n$ = $p$ noves e $n$
zeros:}
\formula{$x = \dfrac{x\cdot10^{n+p} - x\cdot10^{n}}{10^{n+p} - 10^{n}}$}

\SH{$1{,}5\overline{3}$ \quad $(n=1,\,p=1,\,\text{den.}=90)$}
\exemp{%
  $\dfrac{153 - 15}{90} = \dfrac{138}{90}$\\
  MDC$(138,90) = 6 \;\Rightarrow\; \hlm{\dfrac{23}{15}}$%
}

\SH{$2{,}74\overline{63}$ \quad $(n=2,\,p=2,\,\text{den.}=9900)$}
\exemp{%
  $\dfrac{27463 - 274}{9900} = \dfrac{27189}{9900}$\\
  MDC$(27189,9900) = 9 \;\Rightarrow\; \hlm{\dfrac{3021}{1100}}$%
}

\SH{$0{,}5\overline{2}$ \quad $(n=1,\,p=1,\,\text{den.}=90)$}
\exemp{%
  $\dfrac{52 - 5}{90} = \dfrac{47}{90}$\\
  MDC$(47,90) = 1 \;\Rightarrow\; \hlm{\dfrac{47}{90}}$ {\footnotesize(irredut.)}%
}

\end{multicols}

\end{document}
